%Master File:lectures.tex
\newcommand{\rb}{\makebox(0,0){\textcolor{red}{$\bullet$}}}
\newcommand{\yb}{\makebox(0,0){\textcolor{yellow}{$\bullet$}}}
\newcommand{\wb}{\makebox(0,0){$\bullet$}}


\lesson{Project}
\vspace{-1cm}
\begin{center}
\noindent\includegraphics[height=90mm]{mad-scientist.jpg}
\end{center}
\keywords{Details, Ideas, Research Methods}
%%%%%%%%%%%%%%%%%%%%%%% Next Slide %%%%%%%%%%%%%%%%%%%%%%%
\renewcommand{\Outline}{%
\begin{slide}
\section[1]{Outline}

\begin{minipage}{8cm}\raggedright
  \begin{enumerate}\squeeze
    \outlineitem{Rules of the Game}{rules}
    \outlineitem{Ideas}{ideas}
    \outlineitem{How to Do a ML Project}{how}
  \end{enumerate}
\end{minipage}\hfill
\begin{minipage}{13cm}
\ifnum\value{outlineitem}=1
  \includegraphics[width=11cm]{project.jpg}
\fi
\ifnum\value{outlineitem}=2
  \includegraphics[width=13cm]{frankenstein_alive.jpg}
\fi
\ifnum\value{outlineitem}=3
  \includegraphics[width=9cm]{book.jpeg}
\fi
\end{minipage}
\end{slide}
\addtocounter{outlineitem}{1}
}

\setcounter{outlineitem}{1}

%%%%%%%%%%%%%%%%%%%%%%% Next Slide %%%%%%%%%%%%%%%%%%%%%%%
\Outline % Motivation
\toptarget{firstoutline}
%%%%%%%%%%%%%%%%%%%%%%% Next Slide %%%%%%%%%%%%%%%%%%%%%%%

\begin{slide}
\section{Rules}

\begin{PauseHighLight}
  \begin{itemize}
  \item You will work in teams of 3-5\pause
  \item Next week each team will email me a sheet with their team
    members\pause\ and a team name\pauseb
  \item You will give a brief presentation in week 2 (1 minute talk + 1
    minutes of questions/suggestions)\pause
  \item Write a project brief (1 page) for week 3\pause
  \item In week 8 each team submits a 4 page report on data exploration\pause
  \item In week 11 each team submits a 4 page report on using ML\pause
  \end{itemize}
\end{PauseHighLight}

\end{slide}

%%%%%%%%%%%%%%%%%%%%%%% Next Slide %%%%%%%%%%%%%%%%%%%%%%%

\begin{slide}
\section{Help}

\begin{PauseHighLight}
  \begin{itemize}
  \item There is a team of PhD students, Jon, Xiaohao and myself to help\pause
  \item We are only going to help in the Monday afternoon lecture slot\pause
  \item We will give you ideas and steer you in the right direction\pause
  \item We are available (nearly) every week\pause
  \item To give you useful feedback I've split your report\pause
  \end{itemize}
\end{PauseHighLight}

\end{slide}

%%%%%%%%%%%%%%%%%%%%%%% Next Slide %%%%%%%%%%%%%%%%%%%%%%%

\begin{slide}
\section{Marking}

\begin{PauseHighLight}
  \begin{itemize}
  \item The project is worth 20\% (10\% and 10\%)\pause
  \item The marks are assigned on the group reports\pause
  \item If groups don't work I will split you up\pause
  \item Each group can submit 1 page extra to indicate who should get the
    credit (signed by all group members)\pause
  \item The exam is worth 80\% and is the big differentiator\pause
  \item You also get to give group presentations, just for the fun of it!\pauseb
  \end{itemize}
\end{PauseHighLight}

\end{slide}


%%%%%%%%%%%%%%%%%%%%%%% Next Slide %%%%%%%%%%%%%%%%%%%%%%%
\Outline % Ideas
%%%%%%%%%%%%%%%%%%%%%%% Next Slide %%%%%%%%%%%%%%%%%%%%%%%

\begin{slide}
\section{Subject}

\begin{PauseHighLight}
  \begin{itemize}
  \item There is no restriction, but we award marks for work on
    ``advanced machine learning''\pause
  \item We take a very broad view of what that constitutes---its the
    whole process\pause
  \item We are expecting you to go beyond taking a data set and applying
    a machine learning tool out of the box\pause
  \item You can use any tools and any programming language you like
    (clearly acknowledging its use of course)\pause
  \end{itemize}
\end{PauseHighLight}

\end{slide}

%%%%%%%%%%%%%%%%%%%%%%% Next Slide %%%%%%%%%%%%%%%%%%%%%%%

\begin{slide}
\section{What Should You Do?}

\begin{PauseHighLight}
  \begin{itemize}
  \item Find a question, find a data set and make predictions\pause
  \item There are lots of data sources available, datasets text/images
    on the web\pause\ (decent please)\pauseb
  \item Machine learning challenges
    \begin{itemize}
    \item //pascallin2.ecs.soton.ac.uk/Challenges/
    \item //home.web.cern.ch/about/updates/2014/05/higgs-boson-machine-learning-challenge
    \item //www.kaggle.com/\pause
    \end{itemize}
  \item Beat IBM Watson at Jeopardy\pause 
  \end{itemize}
\end{PauseHighLight}

\end{slide}

%%%%%%%%%%%%%%%%%%%%%%% Next Slide %%%%%%%%%%%%%%%%%%%%%%%

\begin{slide}
\section{Covering My Back}

\begin{PauseHighLight}
  \begin{itemize}
  \item Any projects involving
    \begin{itemize}
    \item Humans\pause
    \item Personal data\pause
    \item Creating new intelligent life forms\pauseb
    \end{itemize}
  \item Must have ethical approval\pause
  \end{itemize}
\end{PauseHighLight}

\end{slide}


%%%%%%%%%%%%%%%%%%%%%%% Next Slide %%%%%%%%%%%%%%%%%%%%%%%
\Outline % How
%%%%%%%%%%%%%%%%%%%%%%% Next Slide %%%%%%%%%%%%%%%%%%%%%%%

\begin{slide}
\section{How to Do a Machine Learning Project?}

\begin{PauseHighLight}
  \begin{itemize}
  \item You are going to learn how to do a machine learning project by
    doing it\pause
  \item I'm asking you not only to do the project, but to find out how
    to do a ML project\pause
  \item The book ``Hands-on Machine Learning with Scikit-Learn, Keras,
    and TensorFlow: Concepts, Tools, and Techniques to Build
    Intelligent Systems'' by Aur\'elien G\'eron takes you through this
    process\pause
  \item In the Appendix there is a check list (also on a web
    page---see course page)\pause
  \end{itemize}
\end{PauseHighLight}

\end{slide}

%%%%%%%%%%%%%%%%%%%%%%% Next Slide %%%%%%%%%%%%%%%%%%%%%%%

\begin{slide}
\section[-1]{Eight Steps}
\begin{PauseHighLight}
  \begin{enumerate}\squeeze
  \item Frame the problem and look at the big picture\pause
  \item Get the data\pause
  \item Explore the data to gain insights\pause
  \item Prepare the data to better expose the underlying data patterns
    to Machine Learning algorithms\pause
  \item Explore many different models and short-list the best ones\pause
  \item Fine-tune your models and combine them into a great solution\pause
  \item Present your solution\pause
  \item Launch, monitor, and maintain your system\pause
  \end{enumerate}
\end{PauseHighLight}
\end{slide}

%%%%%%%%%%%%%%%%%%%%%%% Next Slide %%%%%%%%%%%%%%%%%%%%%%%

\begin{slide}
  \section[-2]{1. Frame the Problem and Look at the Big Picture}

  \begin{PauseHighLight}
  \begin{enumerate}\squeeze
  \item Define the objective in business terms\pause
  \item How will your solution be used?\pause
  \item What are the current solutions/workarounds (if any)?\pause
  \item How should you frame this problem (supervised/unsupervised,
    online/offline, etc.)?\pause
  \item How should performance be measured?\pause
  \item Is the performance measure aligned with the business
    objective?\pause
  \item What would be the minimum performance needed to reach the
    business objective?\pause
  \item What are comparable problems? Can you reuse experience or
    tools?\pause
  \item Is human expertise available?\pause
  \item How would you solve the problem manually?\pause
  \item List the assumptions you (or others) have made so far\pause
  \item Verify assumptions if possible\pause
  \end{enumerate}
\end{PauseHighLight}

\end{slide}

%%%%%%%%%%%%%%%%%%%%%%% Next Slide %%%%%%%%%%%%%%%%%%%%%%%

\begin{slide}
\section[-2]{2. Get the Data}

\begin{PauseHighLight}
  \begin{itemize}\squeeze
  \item \emph{Note: automate as much as possible so you can easily get fresh
      data}\pause
  \end{itemize}
  \begin{enumerate}\squeeze
  \item List the data you need and how much you need\pause
  \item Find and document where you can get that data\pause
  \item Check how much space it will take\pause
  \item Check legal obligations, and get authorization if necessary\pause
  \item Get access authorizations\pause
  \item Create a workspace (with enough storage space)\pause
  \item Get the data\pause
  \item Convert the data to a format you can easily manipulate
    (without changing the data itself)\pause
  \item Ensure sensitive information is deleted or protected (e.g.,
    anonymized)\pause
  \item Check the size and type of data (time series, sample,
    geographical, etc.)\pause
  \item Sample a test set, put it aside, and never look at it (no data
    snooping!)\pause
  \end{enumerate}
\end{PauseHighLight}

\end{slide}

%%%%%%%%%%%%%%%%%%%%%%% Next Slide %%%%%%%%%%%%%%%%%%%%%%%

\begin{slide}
  \section[-2]{3. Explore the Data}

\begin{PauseHighLight}
  \begin{itemize}\squeeze
  \item \emph{Note: try to get insights from a field expert for these
      steps.}\pause
  \end{itemize}
  \begin{enumerate}\squeeze
  \item Create a copy of the data for exploration (sampling it down to
    a manageable size if necessary)\pause
  \item Create a Jupyter notebook to keep a record of your data
    exploration\pause
  \item Study each attribute and its characteristics:
    \begin{itemize}
    \item Name
    \item Type (categorical, int/float, bounded/unbounded, text,
      structured, etc)
    \item Number of missing values
    \item Noisiness and type of noise (stochastic, outliers, rounding
      errors, etc.)\pause
    \item Possibly useful for the task?
    \item Type of distribution (Gaussian, uniform, logarithmic, etc)\pause
    \end{itemize}
  \item For supervised learning tasks, identify the target
    attribute(s)\pause
  \item Visualise the data\pause
  \item Study the correlations between attributes\pause
  \item Study how you would solve the problem manually\pause
  \item Identify the promising transformations you may want to
    apply\pause
  \item Identify extra data that would be useful\pause
  \item Document what you have learned\pause
  \end{enumerate}
\end{PauseHighLight}

\end{slide}

%%%%%%%%%%%%%%%%%%%%%%% Next Slide %%%%%%%%%%%%%%%%%%%%%%%

\begin{slide}
  \section[-2]{4. Prepare the Data}

\begin{PauseHighLight}\squeeze
  \begin{itemize}
  \item Notes:
    \begin{itemize}
    \item Work on copies of the data (keep the original dataset
      intact)\pause
    \item Write functions for all data transformations you apply, for
      five reasons:
      \begin{itemize}
      \item So you can easily prepare the data the next time you get a
        fresh dataset
      \item So you can apply these transformations in future project
      \item To clean and prepare the test set
      \item To clean and prepare new data instances once your solution
        is live
      \item To make it easy to treat your preparation choices as
        hyperparameters\pause
      \end{itemize}
    \end{itemize}
  \end{itemize}
\end{PauseHighLight}


\end{slide}

%%%%%%%%%%%%%%%%%%%%%%% Next Slide %%%%%%%%%%%%%%%%%%%%%%%

\begin{slide}
\section[-2]{4. Prepare the Data continued}

\begin{PauseHighLight}\squeeze
  \begin{enumerate}
  \item Data cleaning
    \begin{itemize}
    \item Fix or remove outliers (optional)
    \item Fill in missing values (e.g., with zero, mean, median…) or
      drop their rows (or columns)\pause
    \end{itemize}
  \item Feature selection (optional):
    \begin{itemize}
    \item Drop the attributes that provide no useful information for
      the task\pause
    \end{itemize}
  \item Feature engineering, where appropriate:
    \begin{itemize}
    \item Discretize continuous features
    \item Decompose features (e.g., categorical, date/time, etc.)
    \item Add promising transformations of features (e.g., $\log(x)$,
      $\sqrt{x}$, $x^2$, $\e{x}$ , etc.)\pause
    \item Aggregate features into promising new features\pause
    \end{itemize}
  \item Feature scaling: standardise or normalise features\pause
  \end{enumerate}
\end{PauseHighLight}


\end{slide}


%%%%%%%%%%%%%%%%%%%%%%% Next Slide %%%%%%%%%%%%%%%%%%%%%%%

\begin{slide}
  \section[-2]{5. Short-List Promising Models}

\begin{PauseHighLight}
  \begin{enumerate}\squeeze
  \item Train many quick and dirty models from different categories
    (e.g., linear, naiveBayes, SVM, Random Forests, neural net, etc.)
    using standard parameters\pause 
  \item Measure and compare their performance
    \begin{itemize}
    \item For each model, use N-fold cross-validation and compute the
      mean and standard deviation of the performance measure on the N
      folds\pause
    \end{itemize}
  \item Analyse the most significant variables for each
    algorithm\pause
  \item Analyse the types of errors the models make
    \begin{itemize}
    \item What data would a human have used to avoid these errors?\pause
    \end{itemize}
  \item Have a quick round of feature selection and engineering\pause
  \item Have one or two more quick iterations of the five previous
    steps\pause
  \item Short-list the top three to five most promising models,
    preferring models that make different types of errors\pause
  \end{enumerate}
\end{PauseHighLight}


\end{slide}

%%%%%%%%%%%%%%%%%%%%%%% Next Slide %%%%%%%%%%%%%%%%%%%%%%%

\begin{slide}
\section[-2]{6. Fine-Tune the System}

\begin{PauseHighLight}\squeeze
  \begin{itemize}
  \item Notes:
    \begin{itemize}
    \item You will want to use as much data as possible for this step,
      especially as you move toward the end of fine-tuning
    \item As always automate what you can\pause
    \end{itemize}
  \end{itemize}
  \begin{enumerate}
  \item Fine-tune the hyperparameters using cross-validation
    \begin{itemize}
    \item Treat your data transformation choices as hyperparameters,
      especially when you are not sure about them (e.g., should I
      replace missing values with zero or with the median value? Or
      just drop the rows?)
    \item Unless there are very few hyperparameter values to explore,
      prefer random search over grid search. If training is very long,
      you may prefer a Bayesian optimisation approach\pause
    \end{itemize}
  \item Try Ensemble methods. Combining your best models will often
    perform better than running them individually\pause
  \item Once you are confident about your final model, measure its
    performance on the test set to estimate the generalization error\pause
  \end{enumerate}
\end{PauseHighLight}

\end{slide}

%%%%%%%%%%%%%%%%%%%%%%% Next Slide %%%%%%%%%%%%%%%%%%%%%%%

\begin{slide}
  \section[-2]{7. Present Your Solution}

\begin{PauseHighLight}\squeeze
  \begin{enumerate}
  \item Document what you have done\pause
  \item Create a nice presentation
    \begin{itemize}
    \item Make sure you highlight the big picture first\pause
    \end{itemize}
  \item Explain why your solution achieves the business
    objective\pause
  \item Don’t forget to present interesting points you noticed along
    the way
    \begin{itemize}
    \item Describe what worked and what did not
    \item List your assumptions and your system’s limitations\pause
    \end{itemize}
  \item Ensure your key findings are communicated through beautiful
    visualizations or easy-to-remember statements (e.g., “the median
    income is the number-one predictor of housing prices”)\pause
  \end{enumerate}
\end{PauseHighLight}
\end{slide}

%%%%%%%%%%%%%%%%%%%%%%% Next Slide %%%%%%%%%%%%%%%%%%%%%%%

\begin{slide}
\section[-2]{8. Launch}

\begin{PauseHighLight}\squeeze
  \begin{itemize}
  \item Get your solution ready for production (plug into production
    data inputs, write unit tests, etc.)\pause
  \item Write monitoring code to check your system’s live performance
    at regular intervals and trigger alerts when it drops
    \begin{itemize}
    \item Beware of slow degradation too: models tend to “rot” as data
      evolves
    \item Measuring performance may require a human pipeline (e.g.,
      via a crowd-sourcing service)
    \item Also monitor your inputs’ quality (e.g., a malfunctioning
      sensor sending random values, or another team’s output
      becoming stale). This is particularly important for online
      learning systems\pause
    \end{itemize}
  \item Retrain your models on a regular basis on fresh data (automate
    as much as possible)\pause
  \end{itemize}
\end{PauseHighLight}

\end{slide}

%%%%%%%%%%%%%%%%%%%%%%% Next Slide %%%%%%%%%%%%%%%%%%%%%%%

\begin{slide}
  \section{In Conclusion}

\begin{PauseHighLight}
  \begin{itemize}
  \item Its a long list\pause
  \item Not everything might be relevant\pause
  \item But there is quite a lot of wisdom, so review the list
    regularly during the project\pause
  \item You are doing this for you!\pause
  \item For me, I only want two short reports\pauseb---I've got to
    mark them\pauseb
  \end{itemize}
\end{PauseHighLight}


\end{slide}
